

Ces
relations étant isobares, le quotient est une @@M0@@-algèbre graduée @@M1@@. On définit un morphisme @@M2@@ par :
@@M3@@
Les relations imposées montrent que c'est une théorie de Chern de @@M4@@ à valeurs dans @@M5@@. Si de plus on a une théorie de Chern @@M6@@, à valeurs dans une @@M7@@-algèbre graduée @@M8@@, on a de façon unique un homomorphisme : @@M9@@, tel que @@M10@@ se factorise par @@M11@@; il est défini en envoyant, pour tout @@M12@@, et tout @@M13@@, l'élément @@M14@@ sur le @@M15@@-ième coefficient de @@M16@@ (ce qui est bien défini car @@M17@@ est un @@M18@@-homomorphisme).

Supposons maintenant @@M19@@ de caractéristique nulle, et @@M20@@ de filtration discrète~; soit @@M21@@ tel que @@M22@@. Alors l'algèbre graduée @@M23@@ est nulle en degrés @@M24@@. Pour le montrer, il faut montrer que pour @@M25@@, et @@M26@@, on a :
@@M27@@
Or par hypothèse, on a :
@@M28@@
En appliquant le @@M29@@-homomorphisme @@M30@@, on obtient :
@@M31@@
Or d'après (6.6.4) et (6.5.3), ceci entraîne :
@@M32@@
Comme @@M33@@ est de caractéristique nulle, on a le résultat voulu.

La propriété universelle de @@M34@@ appliquée à l'homomorphisme canonique @@M35@@ donne un diagramme commutatif :
